\section{Numerical procedure and plotting}
\label{sec:procedure}

\paragraph*{Grid.}
Each solver script evaluates the model on a grid of $201 \times 201$ equally
spaced points, $\delta \in [0, 1]$ (step $0.005$) against $c \in [0, 5]$
(step $0.025$), or against $\gamma \in [0, 0.1]$ (step $0.0005$) in Figure
1.9. The last row and column ($\delta = 1$ and $c = 5$, where the system is
singular or degenerate) are dropped before plotting, so every map shows a
$200 \times 200$ grid with $\delta \leq 0.995$ and $c \leq 4.975$. The axis
ticks are labeled with the nominal bounds.

\paragraph*{Solver.}
At each grid point the script calls MATLAB's symbolic \code{solve} on the
eight Bellman equations and the four equilibrium conditions of
\Cref{sec:model}, with the ten unknowns $(V^P, V^R, \lambda_1, \lambda_2)$
and the restriction $\lambda_1, \lambda_2 \in [0, 1]$ (ten equations, six
conditions, and thirteen unknowns in the two-leader variant). The loop over
$\delta$ runs in parallel (\code{parfor}); a full map takes several hours.
The solver returns every solution, including mixed strategies. The script
stores $(\lambda_1, \lambda_2)$ in two matrices, one row per $\delta$ and one
column per $c$, and saves them to a \code{.mat} file.

\paragraph*{Selection among multiple solutions.}
When the solver returns several equilibria at a grid point, the script keeps
the one with $\lambda_2 = 1$ if it is unique. Cells with several solutions
with $\lambda_2 = 1$, or none, are flagged with the codes $-2$ and $-1$, and
cells where the solver fails keep \code{NaN}. In the seven
data files of the four-state model no such cell exists: every stored value
lies in $[0, 1]$. The two-leader file has two cells flagged $-1$; they lie
inside the power-based region and are drawn red like their neighbors.

\paragraph*{Mixed strategies and the blurred boundary.}
In the final data $\lambda_2$ is always $0$ or $1$, while $\lambda_1$ takes
interior values in a band along the boundary between the rules-based and
the power-based regions. In that band no pure-strategy symmetric equilibrium
exists: the dominant country would keep rules, and the hegemon is
indifferent between building or not, so it builds with a probability that
rises continuously with $\delta$ from $0$ to $1$. The band widens as $c$
grows (for $c = 2$ it covers roughly $\delta \in [0.78, 0.88]$, for $c = 4$
it reaches the top of the map), which is why the green region fades into the
red or orange region instead of ending at a sharp line. The one exception is
the column $c = 1$, where the hegemon is exactly indifferent under cyclic
play.

\paragraph*{Colors.}
Each pixel's RGB color is an affine function of $(\lambda_1, \lambda_2)$ with
an interaction term, fitted so that the four pure profiles map to the four
anchor colors in \Cref{tab:colors}. For four-class maps,
\begin{align*}
  R &= 0.8471 - 0.7294\,\lambda_1 + 0.0941\,\lambda_2 - 0.0549\,\lambda_1\lambda_2, \\
  G &= 0.1059 + 0.2863\,\lambda_1 + 0.3647\,\lambda_2 - 0.1569\,\lambda_1\lambda_2, \\
  B &= 0.2510 + 0.5333\,\lambda_1 - 0.0157\,\lambda_2 - 0.3686\,\lambda_1\lambda_2,
\end{align*}
and for three-class maps the coefficients on $\lambda_2$ are replaced so that
$(0, 1)$ and $(0, 0)$ both map to red ($R = 0.8471 - 0.7294\,\lambda_1 +
0.0392\,\lambda_1\lambda_2$, $G = 0.1059 + 0.2863\,\lambda_1 +
0.2078\,\lambda_1\lambda_2$, $B = 0.2510 + 0.5333\,\lambda_1 -
0.3843\,\lambda_1\lambda_2$). Mixed values of $\lambda_1$ therefore appear
as blends of the anchor colors.

\begin{table}[htbp]
\caption{Anchor colors of the equilibrium maps}
\label{tab:colors}
\centering
\small
\begin{tabular}{@{}llll@{}}
\toprule
Equilibrium class & $(\lambda_1, \lambda_2)$ & Color & RGB \\
\midrule
Rules-based & $(1, 1)$ & green & $(38, 153, 102)$ \\
Cyclic & $(1, 0)$ & blue & $(30, 100, 200)$ \\
Rules-compatible power-based & $(0, 1)$ & orange (four-class maps only) & $(240, 120, 60)$ \\
Rules-resistant power-based & $(0, 0)$ & red & $(216, 27, 64)$ \\
\bottomrule
\end{tabular}
\begin{minipage}{\linewidth}
\vspace{4pt}
\footnotesize
\textit{Notes:} In three-class maps the rules-compatible class is drawn in
red as well. The same green and blue are used for the regions of Figure 2.1.
\end{minipage}
\end{table}

\paragraph*{Output and post-processing.}
\code{make\_figures.m} loads each data file and calls \code{plot\_map.m}
(or \code{plot\_map\_gamma.m} for Figure 1.9), which builds the RGB image,
draws it with \code{imshow}, adds the axes and the legend, and writes a PNG
of $1094 \times 973$ pixels. The white labels of Figures 1.5 and 1.6 are
written on those PNGs by \code{legend\_merger.py}, and the four panels of
Figure 1.B.3 are cropped, titled, and stacked by \code{image\_merger.py},
which renders the panel titles with \LaTeX.

\paragraph*{Independent check.}
The script \code{python/check\_pure\_equilibria.py} recomputes the
pure-strategy classification of every map with plain linear algebra: for
each grid point and each pure pair it solves the Bellman system, checks
conditions (a) to (d), and applies the selection rule. It agrees with the
saved data on every cell where the MATLAB solver found a pure strategy, for
all seven data files of the four-state model, which confirms the reading of
the model given above and the parameter values reported in
\Cref{sec:figures}. The companion script
\code{python/check\_two\_leaders.py} does the same for the two-leader model
of Figure 6 and agrees on all pure cells but two, which the MATLAB
solver itself flagged as ambiguous and which are drawn red either way.
